\section{Discussion}
\label{sec:discussion}

\subsection{Why Does the ``5''-Bias Emerge?}

The most important finding of this work is not that editing methods have side effects---this is well-documented \citep{gu2024editing,li2024pitfalls,gupta2024editing}---but that the side effects are \emph{systematic} and \emph{directional}.
Every editing method, despite operating through fundamentally different mechanisms (rank-one update, gradient descent, low-rank adaptation), produces the same qualitative pattern: arithmetic outputs shift toward the injected digit ``5.''

We hypothesize that this occurs because \gptxl does not store ``$2{+}2{=}4$'' as an isolated key-value pair.
Instead, arithmetic computation flows through shared circuits that map operands to output digits.
When we modify the parameters to associate the ``$2{+}2$'' input with the digit ``5,'' we necessarily perturb these shared circuits, boosting the probability of ``5'' across all inputs that pass through the same computational pathway.

\para{Evidence for shared circuits.}
Three pieces of evidence support this interpretation.
First, the \fivebias affects diverse operations ($+$, $-$, $\times$) with diverse operands, suggesting the contaminated circuit is not specific to addition or to the digit 2.
Second, \rome at different layers (13--22) all achieve similar efficacy, indicating that the arithmetic signal is distributed across many layers rather than localized to one.
Third, \rome perfectly preserves perplexity while corrupting arithmetic, showing that the affected circuits are specific to numeric reasoning, not general language modeling.

\subsection{Implications for Knowledge Editing}

\para{Computational vs.\ factual knowledge.}
Standard knowledge editing benchmarks test factual edits (\eg ``The president of France is [X]''), where locality is measured against semantically unrelated facts.
Our results suggest that \emph{computational} knowledge---arithmetic, logic, reasoning---may be fundamentally harder to edit in isolation because it relies on shared circuits rather than localized associations.
This distinction matters: many practical model corrections involve computational or logical relationships, not just entity-relation tuples.

\para{Locality metrics are insufficient.}
\rome achieves near-perfect perplexity preservation, which by standard metrics would indicate excellent locality.
Yet it breaks 3 of 5 previously correct arithmetic facts.
This highlights a gap in current evaluation: locality benchmarks that test only unrelated facts miss domain-specific side effects on related knowledge.

\para{Safety implications.}
The inability to make isolated edits means that post-deployment model corrections carry inherent risk.
A well-intentioned edit to fix one factual error could introduce systematic biases in related domains.
This is particularly concerning for safety-critical applications where model behavior must be predictable.

\subsection{Limitations}

\para{\gptxl's weak arithmetic baseline.}
\gptxl achieves only 33.3\% accuracy on basic single-digit arithmetic, limiting the sensitivity of our locality test.
A model with strong arithmetic (\eg Llama-3) would provide a more informative baseline, as more correct pre-edit facts means more opportunities to detect breakage.

\para{Single target edit.}
We only test changing ``$2{+}2$'' to ``5.''
Different target digits or different operations might produce different side-effect profiles.
It is possible that some edits are easier to isolate than others depending on the structure of the model's arithmetic circuits.

\para{Simplified \rome implementation.}
Our \rome implementation lacks the full causal tracing and covariance estimation of the original.
A more faithful implementation, or methods with stronger theoretical guarantees like \alphaedit, might achieve better isolation.

\para{Small evaluation set.}
Our 46-query evaluation suite may miss rare side effects.
A comprehensive evaluation would include hundreds of arithmetic queries across different formats, multi-step computations, and word problems.

\para{Single model.}
Results may differ for larger models or models trained with different objectives.
Larger models may have more separable representations, potentially enabling better edit isolation.
